\section{Conclusion}
\label{sec:conclusion}
This paper presented a diagnostic protocol for sparse-concept and calibration evaluation in Knowledge Tracing: train-only KC-frequency stratification with an explicit strict cold-start group, learner-based and temporal splits, per-stratum calibration analysis (ECE and the Brier decomposition, with reliability diagrams), sample-size-aware reliability flags, and a reproducible audit workflow (L1--L8) that operationalizes standard safeguards plus a predictive-sanity check. The protocol is intended as reusable evaluation infrastructure for KT research, including graph-based and self-supervised models whose claimed gains concentrate on sparse KCs, where evaluation is least stable. Section~\ref{sec:threshold_simulation} is a single-fold illustration of that use, not a new architecture paper.

The instantiation on three datasets supports a bounded scientific story. Sparse training evidence does not universally degrade discrimination, but it can expose dataset-dependent calibration vulnerability. We characterize the conditions associated with this vulnerability---sparse mass, evaluation support, and frequency--difficulty coupling---and provide reproducible diagnostics for identifying when it matters. On ASSISTments 2012, SimpleKT ECE rises from $0.114$ on dense concepts to $0.154$ on medium concepts and $0.228$ on sparse concepts (Limited flag, $N = 415$). Other datasets show flat or inverted stratum patterns. On the seed-42 learner-based fold of the same dataset, a global gate at $\tau=0.7$ raises SimpleKT false mastery among advances from $0.196$ (dense) to $0.268$ (sparse); that $\Delta\mathrm{FM}$ is positive on all five training seeds (mean $0.047$; four unique partitions). A train-only GKT graph shrinks the seed-42 gap, and a CL4KT adapter is intermediate. This is a simulated decision, not a classroom intervention. A controlled downsampling of training rows for the same originally dense KCs produces a within-KC ECE increase only in some dataset--model cells (Section~\ref{sec:a9_sparsify}), so observational frequency--calibration gradients should not be read as a universal frequency dose--response. Sparse-concept diagnostics are most informative when a dataset contains a meaningful low-frequency tail, those strata have sufficient evaluation support, and calibration reliability varies with training evidence or deployment involves concept cold-start (Section~\ref{sec:a8_when_needed}). Occupancy reporting remains useful even when that empirical priority is low. Claim intensity is therefore conditional: not every KT log should be narrated as a sparse-calibration failure. Finally, channel L8 showed empirical value: it caught a prediction--label alignment artifact that would otherwise have been misread as genuine cold-start failure, and a controlled fault-injection validation later measured sensitivity to additional alignment-fault classes (Appendix~\ref{app:l8_fault}). L8 is reported as sensitive to the classes we tested, not as a detector of every evaluation bug. L1--L7 remain standard safeguards made inspectable; they are not offered as a new auditing methodology.

Future work can extend this diagnostic foundation by evaluating across multiple temporal seeds or cutoffs to capture variance, by incorporating a classical baseline that exploits item difficulty for unseen learners under learner-based splits, by applying post-hoc or stratum-aware calibration (a separate methods question), and by repeating the GKT/CL4KT occupancy-and-gate checks on additional datasets and official checkpoints.
